% =====================================================================
%  CMPT 981 Project - Proposal (NeurIPS 2022 format)
%
%  Build (from this folder):  latexmk proposal_2_azaan.tex
%  latexmkrc puts the PDF in ../docs/ and keeps build files here.
%  References live in ref.bib.
% =====================================================================

\documentclass{article}


% if you need to pass options to natbib, use, e.g.:
%     \PassOptionsToPackage{numbers, compress}{natbib}
% before loading neurips_2022


% ready for submission
%\usepackage{neurips_2022}


% to compile a preprint version, e.g., for submission to arXiv, add add the
% [preprint] option:
%     \usepackage[preprint]{neurips_2022}


% to compile a camera-ready version, add the [final] option, e.g.:
     \usepackage[final]{neurips_2022}


% to avoid loading the natbib package, add option nonatbib:
%    \usepackage[nonatbib]{neurips_2022}


\usepackage[utf8]{inputenc} % allow utf-8 input
\usepackage[T1]{fontenc}    % use 8-bit T1 fonts
\usepackage{hyperref}       % hyperlinks
\usepackage{url}            % simple URL typesetting
\usepackage{booktabs}       % professional-quality tables
\usepackage{amsfonts}       % blackboard math symbols
\usepackage{nicefrac}       % compact symbols for 1/2, etc.
\usepackage{microtype}      % microtypography
\usepackage{xcolor}         % colors

% --- Added for this proposal ------------------------------------------
\usepackage{amsmath, amssymb, amsthm}
\usepackage{mathtools}
\usepackage{bm}
\usepackage{graphicx}
\usepackage{subcaption}
\usepackage{enumitem}
\usepackage[capitalize,noabbrev]{cleveref} % must load after hyperref

% Theorem environments
\newtheorem{theorem}{Theorem}
\newtheorem{proposition}[theorem]{Proposition}
\newtheorem{lemma}[theorem]{Lemma}
\theoremstyle{definition}
\newtheorem{definition}[theorem]{Definition}

% Handy macros
\newcommand{\R}{\mathbb{R}}
\newcommand{\E}{\mathbb{E}}
\newcommand{\1}{\mathbf{1}}
\newcommand{\KL}{\mathrm{KL}}
\newcommand{\pref}{p_{\mathrm{ref}}}
\newcommand{\piref}{\pi_{\mathrm{ref}}}
\DeclareMathOperator{\logit}{logit}
\DeclareMathOperator{\Bin}{Bin}
\DeclareMathOperator*{\argmax}{arg\,max}
\newcommand{\norm}[1]{\left\lVert #1 \right\rVert}

% Inline TODO notes; comment out the first line and uncomment the
% second to hide them all before submitting.
\newcommand{\todo}[1]{\textcolor{red}{[TODO: #1]}}
% \newcommand{\todo}[1]{}


\title{GRPO unraveled: A Small-Scale Empirical Study of RLVR\\
  CMPT 981 Project Proposal}


\author{%
  Ankush Singh\\
  Simon Fraser University\\
  \texttt{asa715@sfu.ca} \\
  \And
  Simon Doan\\
  Simon Fraser University\\
  \texttt{vtd@sfu.ca} \\
  \And
  Azaan Rehman\\
  Simon Fraser University\\
  \texttt{ara199@sfu.ca} \\
}


\begin{document}


\maketitle



\begin{abstract}
  Reinforcement learning with verifiable rewards (RLVR), and in particular
  Group Relative Policy Optimization (GRPO), is now the standard recipe for
  training reasoning in large language models. Recent theory makes concrete
  predictions about its dynamics: how many sampled groups carry a learning
  signal, whether GRPO amplifies the success probability, and when a step
  size is too large. These predictions are derived for exact gradients and
  idealized policies, and are hard to check at LLM scale, where every run is
  expensive and the true success probability is unknown. We propose an
  experiment-centred project that tests them on a small synthetic task with
  a binary verifier. Its difficulty is set exactly by the answer length, and
  the policy's success probability can be computed in closed form, so runs
  take seconds on a CPU. We vary the initial success rate, the number of
  sampled answers per prompt, reward sparsity, KL strength, and step size,
  and compare each result with the corresponding theoretical prediction. As
  a small theoretical contribution, we extend the success-amplification
  fixed point of \citet{mroueh2025rlvr} from infinite to finite group size,
  which gives closed forms for mean-only and leave-one-out baselines and a
  testable prediction of how the group size shifts GRPO's fixed point. We
  also show that on-policy SGD shares these fixed points with Mroueh's
  exact iteration but not their stability. We close with a confirmation on
  a 0.5B-parameter LLM.
\end{abstract}

% ---------------------------------------------------------------------
\section{Introduction and Motivation}
\label{sec:intro}

Reinforcement learning (RL) learns a policy of actions from a reward signal
that is often sparse and delayed. In the context of large language models
(LLMs), RL is used to align the model's responses with human preferences or
task-specific objectives, most notably through reinforcement learning from
human feedback (RLHF)~\citep{ouyang2022instructgpt}. However, RLHF can be
expensive and time-consuming, as it requires humans to label which responses
they prefer, and those labels are used to train a reward model.

In 2025, DeepSeek-R1-Zero~\citep{guo2025deepseekr1} showed that a language
model could learn long chain-of-thought reasoning from a rule-based reward,
mainly a binary check of whether the final answer is correct, without any
human feedback. This approach is called reinforcement learning with
verifiable rewards (RLVR). The model was trained using Group Relative Policy
Optimization (GRPO)~\citep{shao2024deepseekmath}, which samples multiple
answers per prompt and normalizes the advantages by the group standard
deviation. Since then, practitioners have relied on rules of thumb, such as
sampling more answers for problems with a low success rate, normalizing
advantages by the group metrics, and averaging the loss over tokens.
However, each of these choices changes the fundamental optimization problem,
and it is still unclear which ones matter and why.

Initial theoretical work provides precise answers in simplified
settings~\citep{mroueh2025rlvr,suk2025gradientgap,thrampoulidis2026passk,bay2026grpostd}.
Examples include the success amplification property of
GRPO~\citep{mroueh2025rlvr}, which states that, with a KL penalty, the
success probability settles at a fixed point higher than that of the
reference model, and the gradient gap~\citep{suk2025gradientgap}, the
difference between the expected score of correct and incorrect responses,
which is used to derive a step-size threshold below which training
converges. However, these results are derived under simplified conditions,
far from the stochastic, finite-sample regime in which GRPO is actually run.
For example, they assume exact expectations and, in the case of
\citet{mroueh2025rlvr}, an exactly solved inner optimization problem.
Testing them directly is also hard, because LLM experiments cannot isolate a
single variable: the success rate, task difficulty, and response length all
change together.

Our project aims to bridge the gap between theory and practice. We plan
controlled experiments that will allow us to sweep a grid of parameters,
including the initial success rate, group size, reward sparsity, KL
strength, and step size. By changing one parameter at a time and observing
the effect on the learning dynamics, we can test the theoretical predictions
in a setting closer to practice, with finite samples, stochastic gradients,
and a learned policy. These results will offer further insight into the
dynamics of RLVR and GRPO in practical settings. This can lead to a better
understanding of the optimization problem and help guide how these settings
are chosen when training LLMs with verifiable rewards.

% ---------------------------------------------------------------------
\section{Related Work}
\label{sec:related}

\paragraph{GRPO and RLVR.}
GRPO~\citep{shao2024deepseekmath} drops the value network of Proximal Policy
Optimization (PPO). For each prompt, it samples a group of $G$ answers and uses the
group-normalized reward $(r_i - \bar r)/\mathrm{std}(r)$ as the advantage.
DeepSeek-R1~\citep{guo2025deepseekr1} scaled this up with rule-based verifiers,
$G=16$, and a small KL coefficient that keeps the policy close to a reference model.
A later study, Dr.~GRPO~\citep{liu2025drgrpo}, argued that the std-normalization and
the per-response length normalization introduce biases. RLOO~\citep{ahmadian2024rloo}
instead uses a leave-one-out baseline, the mean reward of the other $G-1$ answers,
which reduces variance without biasing the gradient. We use all three as advantage
variants in our experiments to see how they affect the learning dynamics.

\paragraph{Success amplification~\citep{mroueh2025rlvr}.}
In 2025, Mroueh analyzed GRPO in the RLVR setting and identified its success
amplification property. The paper shows that, for binary rewards with a reverse-KL
penalty and no clipping, GRPO's objective can be seen as a contrastive loss between
correct and incorrect samples taken from the previous policy. Solving the
KL-regularized problem exactly at each iteration shows that the success probability
follows a recurrence $p_n = h(p_{n-1})$ for an explicit function $h$. Every fixed point
of this recurrence satisfies $p^\ast > \pref$, which means that GRPO increases the
success probability above that of the reference model. In this project, we will test
whether SGD training reaches the predicted $p^\ast$ for different values of $\beta$
(the strength of the KL penalty) and for finite $G$. We will also compare this
reverse-KL penalty, $\KL(\pi_\theta \,\|\, \piref)$, with the forward-KL penalty
$\KL(\piref \,\|\, \pi_\theta)$, which will cover the difference between mode-seeking
and mass-covering policies.

\paragraph{Gradient gap and step-size thresholds~\citep{suk2025gradientgap}.}
The paper defines the gradient gap, the difference between the expected
score of correct and incorrect responses, and proves a sufficient step-size
condition for convergence together with lower-bound constructions
where a constant step makes performance decrease to zero. With exact
gradients, the threshold scales as $1/T$ in the response length $T$ and
shrinks with the success rate through $(1-J)$; the paper argues that GRPO's
$1/|y|$ normalization provides exactly this scaling and Dr.~GRPO's does not.
Its experiments use a single step size. We test whether a sharp transition
appears in a stochastic, learned-policy setting, and how it moves with
response length when responses of different lengths share one learning
rate. \citet{ge2026normalization} sharpen this question: they show that
std-normalization acts as an adaptive step size matched to local
curvature, with a speed-up that depends on how much the per-prompt reward
variance differs across prompts.

\paragraph{Group size and informative prompts.}
A group with all-equal rewards has zero advantage, and a common practice is
to filter out such groups. \citet{bay2026grpostd} show that GRPO, Dr.~GRPO
and DAPO differ only in how they use the group standard deviation, derive a
closed-form group-size law in $G\,p(1-p)$, and check it on a Bernoulli toy
model at $G=8$ and by subsampling logged rollouts.
\citet{nie2026starvation} prove that when the success rate varies across
prompts, the fraction of all-equal groups exceeds the prediction from the
average success rate, and that a fixed-reference advantage $A = 2r - 1$
outperforms Dr.~GRPO at $G = 4$ on GSM8K. \citet{plyusov2026fgrpo} show that
a group can have mixed rewards yet miss rare correct solutions, with a
non-monotonic dependence on $G$.

We do not claim the informative-group fraction (Eq.~\eqref{eq:informative})
or the $p(1-p)$ scaling as new. The closest prior work to ours is
\citet{bay2026grpostd}, \citet{nie2026starvation} and
\citet{plyusov2026fgrpo}, each of which varies $G$ in theory or in a toy
model. Our contribution is a controlled \emph{training} grid over $G$,
$p_0$, KL strength and step size, with every quantity measured exactly,
compared against these closed forms.

\paragraph{Pass@$K$ advantage shaping.}
\citet{thrampoulidis2026passk} show that advantage-shaping
heuristics, such as hard-example upweighting, implicitly maximize surrogate
rewards such as pass@$K$, and derive an unbiased leave-one-out pass@$K$
advantage. We test whether optimizing pass@$K$ preserves response diversity
and at what cost to pass@1.

\paragraph{Entropy and diversity collapse.}
\citet{yue2025doesrl} show that RLVR-trained models beat their base model
at small $k$ but lose to it at large $k$, which suggests that RLVR mostly
sharpens what the base model can already do.

% ---------------------------------------------------------------------
\section{Problem Formulation}
\label{sec:problem}

We start with a prompt $x$ and a policy $\pi_\theta(y \mid x)$ that generates
a response $y$. A verifier returns a reward $r(x, y) \in \{0, 1\}$, where
the policy has a chance of success $p(x)$ for each prompt. We sample a group
of $G$ responses per prompt, compute advantages $A_i$ within the group,
and take a gradient step on $$J(θ) = E[r] − β·KL(πθ ‖ πref)$$.

We define the following variables:
Initial success rate: p₀, Group size:G, Reward sparsity	set by answer
length L and vocabulary K: p₀ = K^(−L), KL strength	β, Step size	η, Advantage
type:	GRPO, mean-only, RLOO, pass@k.

Our policy also covers few variants. Bandit policy: L = 1, so the policy is
a single softmax over K options. Synthetic task: L>1, so the policy is a
product of L softmaxes. Small LLM, so the policy is a real world example.

% ---------------------------------------------------------------------
\section{Plan}
\label{sec:experiments}

Our main goal is to test the limits of theoretical predictions mentioned in
related work. Those works assume simplified conditions, we see where they
break when those conditions are not fullfilled.

We will run a grid of experiments over the parameters mentioned in the
Problem Formulation section. Ideally we will be able to cover the entire grid,
but we prioritize finding the limits of the theoretical methods first, and
then properly cover the strcture of underlying functions around those limits.

We defin following core experiments in order:
\begin{description}[leftmargin=2.2em, labelindent=0em]
  \item[E1] \textbf{Initial success rate $\times$ group size.} Our first test
    is to check the affects of initial success rate and group size on the
    learning speed. We will also compare our measured learning speeds with
    the group-size law of \citet{bay2026grpostd}.

  \item[E2] \textbf{Success amplification~\citep{mroueh2025rlvr}.} With a
    simpler setting of $L = 1$, we will test the success amplification property
    of GRPO. We will compare the measured fixed points with finite group size
    compared to the infinite group size fixed points of \citet{mroueh2025rlvr}.

  \item[E3] \textbf{Step-size threshold~\citep{suk2025gradientgap}.} We will test
    the step-size threshold property of GRPO. We will sweep the learning rate
    over a log grid and check for a narrow learning-rate range where success
    first rises and then stalls or collapses. We will also check if the
    threshold shrinks roughly as $1/L$ and as $p_0$ grows.

  \item[E4] \textbf{Pass@$k$ reward diversity~\citep{thrampoulidis2026passk}.}
    On the multi-solution tasks, we will compare the advantage of pass@$k$.

  \item[E5] \textbf{Escaping very sparse rewards.} On the boundaries of sparse rewards,
    we will test what helps the model escape the collapse.

  \item[E6] \textbf{Small-LLM check.} We will run a small LLM to check the
    theoretical predictions and new insights from the previous experiments.
\end{description}

\paragraph{Current status.}
We have derived \cref{prop:finiteG} and computed the predicted fixed
points and their stability (\cref{tab:omega}, Eq.~\eqref{eq:flow}). The
training harness (task, policy, advantage estimators, and exact metrics)
is not yet implemented; it is the first milestone.

\paragraph{Timeline and roles.}
Weeks 1--2: build the harness and unit-test the exact metrics against
Eq.~\eqref{eq:informative} and the closed-form fixed points; complete E1. Weeks 3--4
(milestone): complete E2 and E3, and set up the E6 pilot run. Weeks 5--6:
extensions E4--E6, ablations, and the final report. Each member owns one
experiment end-to-end (implementation, runs, figures, and the related
section of the report), and all four share the code and the literature
review.

\paragraph{What we expect to learn.}
The theory is derived for exact gradients and simple policies. Where our
curves deviate from it, we will identify the cause: finite $G$ (which
\cref{prop:finiteG} quantifies for E2), stochastic batches, function
approximation shared across prompts, or the factorized multi-token policy.
These deviations are the most useful output of the project.

% ---------------------------------------------------------------------
\bibliographystyle{plainnat}
\bibliography{ref}


% ---------------------------------------------------------------------
\appendix

\section{Proof of \texorpdfstring{\cref{prop:finiteG}}{the Proposition}}
\label{app:proof}

Fix a prompt and write $\pi = \pi_\theta(\cdot \mid x)$ and
$p = \sum_{y : r(y) = 1} \pi(y)$. The advantages are computed from
sampled rewards and are not differentiated, so the expected update
direction of the reward term is
$g = \E\bigl[\tfrac{1}{G}\sum_{i=1}^{G} A_i \nabla_\theta \log \pi(y_i)\bigr]$.
The group is exchangeable, so $g = \E[A_1 \nabla_\theta \log \pi(y_1)]$.
Conditional on $y_1$, the other responses are independent of $y_1$, so
$\E[A_1 \mid y_1]$ depends on $y_1$ only through $r(y_1)$, and equals
$w(1) = a^+_G(p)$ or $w(0) = -a^-_G(p)$. Hence
\[
  g = \sum_y \pi(y)\, w(r(y))\, \nabla_\theta \log \pi(y)
    = \nabla_\theta \sum_y \pi_\theta(y)\, w(r(y)) \Big|_{w \text{ fixed}} .
\]
Adding the exact KL penalty, a stationary point is a stationary point of
$F(\pi) = \sum_y \pi(y) w(r(y)) - \beta\, \KL(\pi \,\|\, \piref)$ at fixed
$w$. When the policy class contains every distribution with full support
(a softmax over all responses), $\nabla_\theta F = 0$ if and only if
$w(r(y)) - \beta\bigl(\log\tfrac{\pi(y)}{\piref(y)} + 1\bigr)$ is constant
in $y$, i.e.\ $\pi(y) \propto \piref(y) \exp(w(r(y))/\beta)$. Summing
over correct and incorrect responses,
\[
  \frac{p}{1-p} = \frac{\pref}{1-\pref}
    \exp\!\Bigl(\frac{a^+_G(p) + a^-_G(p)}{\beta}\Bigr),
\]
which is Eq.~\eqref{eq:fixedpoint}, with $p$ evaluated at the stationary
point itself.

For \texttt{mean-only}, $A_1 = r_1 - \frac{1}{G}\sum_i r_i$, so
$a^+_G = 1 - \frac{1 + (G-1)p}{G} = \frac{G-1}{G}(1-p)$ and
$a^-_G = \frac{G-1}{G}p$, giving $\Omega_G = (G-1)/G$. For \texttt{rloo},
$A_1 = r_1 - \frac{1}{G-1}\sum_{i \ge 2} r_i$, so $a^+_G = 1-p$,
$a^-_G = p$ and $\Omega_G = 1$. For \texttt{grpo} with $G = 2$, the only
informative group has rewards $\{0, 1\}$ and advantages
$\pm \tfrac{1}{2}/\sqrt{\tfrac14 + \epsilon} \approx \pm 1$, so
$a^+_2 \approx 1-p$, $a^-_2 \approx p$ and $\Omega_2 \approx 1$. As
$G \to \infty$, $j/G \to p$ in probability, so
$a^+_G \to (1-p)/\sqrt{p(1-p)+\epsilon}$ and
$a^-_G \to p/\sqrt{p(1-p)+\epsilon}$ by bounded convergence, and
$\Omega_G \to (p(1-p)+\epsilon)^{-1/2}$, the weight in the recurrence of
\citet{mroueh2025rlvr}. If the implementation uses the Bessel-corrected
standard deviation, replace $\frac{c}{G}(1-\frac{c}{G})$ in
Eq.~\eqref{eq:omegaG} by $\frac{c(G-c)}{G(G-1)}$. \hfill$\square$

\paragraph{Derivation of Eq.~\eqref{eq:flow}.}
Take $L = 1$, $K$ answers with one correct answer $c$, and logits
$\ell_k$, so $\pi_k \propto e^{\ell_k}$. Write
$g_k = w(r_k) - \beta(\log\frac{\pi_k}{\pi_{\mathrm{ref},k}} + 1)$ and
$\bar g = \sum_k \pi_k g_k$. The expected-gradient flow on $F$ is
$\dot\ell_k = \pi_k (g_k - \bar g)$. With a zero-initialized head, the
$K-1$ wrong answers start with equal logits and receive equal updates,
so they stay equal, each with probability
$\pi_w = (1-p)/(K-1)$. Let $\Delta = g_c - g_w$. Then
$\dot\ell_c = p(1-p)\Delta$ and $\dot\ell_w = -p\,\pi_w \Delta$. Using
$\pi_{\mathrm{ref},w} = (1-\pref)/(K-1)$,
$\Delta = a^+_G + a^-_G - \beta(z - \logit \pref)$, where
$z = \logit p = \ell_c - \ell_w - \log(K-1)$. Hence
$\dot z = \dot\ell_c - \dot\ell_w = \kappa\,\Delta$ with
$\kappa = p(1-p) + p\,\pi_w > 0$, which is Eq.~\eqref{eq:flow}. A fixed
point is stable when $\partial\Delta/\partial z < 0$, i.e.\ when
$\Omega_G'(p^\ast)\,p^\ast(1-p^\ast) < \beta$. Mroueh's map
$z_n = \logit\pref + \Omega(\sigma(z_{n-1}))/\beta$ has slope
$s = \Omega'(p^\ast)\,p^\ast(1-p^\ast)/\beta$ at a fixed point and is
stable when $|s| < 1$. The two conditions agree when $s > 0$; when
$s < -1$, the map oscillates but the flow is still stable.

\section{Experimental Details}
\label{app:details}

\begin{table}[h]
  \centering
  \begin{tabular}{lp{0.66\linewidth}}
    \toprule
    Component & Setting \\
    \midrule
    Prompt features & $x \sim \mathcal{N}(0, I_{16})$; 512 train / 256 held-out test prompts \\
    Task & $K = 4$ (E1, E3--E5), $K \in \{8, 1000\}$ (E2); $L \in \{1, \dots, 8\}$; Gaussian $W_t$ \\
    Mixed-length task (E3) & $L(x)$ uniform on $\{2, \dots, 8\}$, fixed per prompt \\
    Multi-solution task (E4) & $m$ of $K$ tokens correct per position; unequal initial probabilities on correct tokens \\
    Policy & MLP (2 hidden layers, width 128) on $(x, \text{one-hot } t)$; zero-init head, so $\piref$ is uniform. E2 and E3 also use a tabular softmax per prompt \\
    Optimizer & Adam (E1, E4, E5); plain SGD for E2 and the learning-rate sweep (E3) \\
    Batch & $B = 32$ prompts $\times$ $G$ samples; fixed-compute runs use $B \times G = 256$ \\
    Advantage std & $\sqrt{\bar r(1-\bar r) + \epsilon}$, $\epsilon = 10^{-5}$, no Bessel correction \\
    Seeds & 3 per configuration \\
    Compute & CPU only; each run takes seconds to minutes \\
    \midrule
    LLM (E6) & Qwen2.5-0.5B-Instruct; GSM8K (MATH levels 1--5 as fallback) \\
    LLM training & TRL GRPOTrainer + vLLM; $G \le 16$, lr $5 \times 10^{-6}$, $\beta = 0.1$, 200-token context, temperature 0.1 \\
    LLM compute & H100 MIG slices (40\,GB) on DRAC Fir; $\approx$3 GPU-days budget \\
    \bottomrule
  \end{tabular}
\end{table}

\end{document}
